\cleardoublepage % memastikan bab baru dimulai di halaman ganjil (kanan)
\chapter{ORGANISASI ATAU LINGKUNGAN KERJA PRAKTIK}
\label{chap:organisasi_lingkungan}
\section{Profil Perusahaan}
\label{sec:profil_perusahaan}
Meskipun telah dibahas di dalam Bab \ref{chap:pendahuluan}, profil perusahaan tempat pelaksanaan kerja praktik perlu dijelaskan secara lebih rinci. Tuliskan profil singkat perusahaan tempat pelaksanaan kerja praktik, termasuk sejarah singkat, bidang usaha, visi dan misi perusahaan, serta produk atau layanan utama yang ditawarkan oleh perusahaan tersebut.
Deskripsikan juga posisi perusahaan dalam industri atau pasar yang relevan, serta pencapaian atau penghargaan penting yang telah diraih oleh perusahaan.
Informasi ini dapat diperoleh dari situs web resmi perusahaan, brosur perusahaan, atau sumber informasi lainnya yang terpercaya.
\section{Struktur Organisasi}
\label{sec:struktur_organisasi}
Perusahaan yang dibahas di Subbab \ref{sec:profil_perusahaan} memiliki struktur organisasi. Jelaskan struktur organisasi perusahaan tempat pelaksanaan kerja praktik.
Sertakan bagan struktur organisasi yang menggambarkan hierarki dan hubungan antar departemen atau divisi dalam perusahaan.
Deskripsikan juga peran dan tanggung jawab masing-masing departemen atau divisi yang relevan dengan pelaksanaan kerja praktik.
Jika memungkinkan, jelaskan juga posisi penulis dalam struktur organisasi perusahaan selama pelaksanaan kerja praktik.
\section{Unit Kerja}
\label{sec:unit_kerja}
Deskripsikan unit kerja atau departemen tempat penulis melaksanakan kerja praktik, termasuk fungsi dan tugas utama unit kerja tersebut dalam perusahaan.
Jelaskan juga bagaimana unit kerja tersebut berkontribusi terhadap tujuan dan operasional perusahaan secara keseluruhan.
Jika ada, sertakan informasi tentang proyek atau inisiatif khusus yang sedang dijalankan oleh unit kerja selama periode kerja praktik.
\section{Lingkungan Kerja} 
\label{sec:lingkungan_kerja}
Jelaskan lingkungan kerja di perusahaan tempat pelaksanaan kerja praktik, termasuk budaya kerja, suasana kerja, dan interaksi antar karyawan.
Deskripsikan juga fasilitas dan sumber daya yang tersedia bagi karyawan, seperti ruang kerja, peralatan, dan teknologi yang digunakan dalam operasional sehari-hari.
Jika relevan, jelaskan juga tantangan atau peluang yang dihadapi dalam lingkungan kerja tersebut, serta bagaimana penulis menyesuaikan diri dengan lingkungan kerja selama pelaksanaan kerja praktik.








